\section{2013 Applied \& Computational Mathematics}

\subsection{Individual}

\begin{exercise}
%(5 pts) We consider the wave equation $u_{tt}=\Delta u$ in $\mathbb{R}^{3}\times\mathbb{R}_{+}$.

\begin{enumerate}
\item[(a)] A right going pulse with speed 1,
\[
u(x,y,z,t)=1\ \text{for}\ t<x<t+1;\quad u(x,y,z,t)=0\ \text{else},
\]
is clearly a solution to the wave equation. However, it is a discontinuous solution. Explain in which sense it is a solution to the equation.
\item[(b)] %(5 pts) Surprisingly, one can construct smooth progressive wave solutions with speed larger than 1 (superluminal wave). Try a solution of the form
\[
u(x,y,z,t)=v\left(\frac{x-ct}{\sqrt{c^{2}-1}},y,z\right),\quad c\in\mathbb{R}^{3},\ |c|>1.
\]
Derive an equation for $\mathbf{v}$ and show there is a nontrivial solution with compact support in $(y,z)$ for any fixed $x,t$.
\item[(c)] %(5 pts) For any $R>0$, $0<t<R$, show that energy
\[
E(t):=\int_{|\vec{x}|\leq R-t}\left(|u_{t}(\cdot,t)|^{2}+|\nabla u(\cdot,t)|^{2}\right)d\vec{x}
\]
is a decreasing function.
\item[(d)] %(10 pts) Show that smooth superluminal progressive wave solutions of the form
\[
u(\vec{x},t)=v(\vec{x}-\vec{c}t),\quad\vec{c}\in\mathbb{R}^{3},\ |\vec{c}|>1,
\]
cannot have finite energy.

Hint: Using (c) and look at the energy of the solution in various balls.
\end{enumerate}
\end{exercise}

\begin{solution}
\begin{enumerate}
    \item[(a)] The pulse is a \textbf{weak solution} (or distributional solution). A function $u \in L^1_{loc}(\mathbb{R}^3 \times \mathbb{R}_+)$ is a weak solution to the wave equation $u_{tt} - \Delta u = 0$ if for all $\phi \in C_c^\infty(\mathbb{R}^3 \times \mathbb{R}_+)$, we have
    \[
    \iint_{\mathbb{R}^3 \times \mathbb{R}_+} u (\phi_{tt} - \Delta \phi) \, dx \, dt = 0.
    \]
    Since $u(x,y,z,t) = f(x-t)$ with $f = \mathbb{1}_{(0,1)}$, and $f$ is a distribution satisfying $f'' - f'' = 0$, $u$ satisfies the equation in the sense of distributions.
    
    \item[(b)] Let $\xi = \frac{x-ct}{\sqrt{c^2-1}}$. Then $u_t = v_\xi \frac{-c}{\sqrt{c^2-1}}$ and $u_{tt} = v_{\xi\xi} \frac{c^2}{c^2-1}$.
    Also $u_x = v_\xi \frac{1}{\sqrt{c^2-1}}$ and $u_{xx} = v_{\xi\xi} \frac{1}{c^2-1}$.
    The wave equation $u_{tt} = u_{xx} + u_{yy} + u_{zz}$ becomes:
    \[
    \frac{c^2}{c^2-1} v_{\xi\xi} = \frac{1}{c^2-1} v_{\xi\xi} + v_{yy} + v_{zz} \implies v_{\xi\xi} = v_{yy} + v_{zz}.
    \]
    This is the \textbf{2D wave equation} where $\xi$ plays the role of "time". By the property of the wave equation in 2D (or Huygens' principle/domain of dependence), if we provide compactly supported "initial data" at $\xi=0$, the solution $v(\xi, y, z)$ will have compact support in $(y,z)$ for any finite $\xi$. Specifically, let $v(0, y, z) = f(y, z)$ and $v_\xi(0, y, z) = g(y, z)$ be $C_c^\infty$ functions. Then for any fixed $\xi$, $v$ is compactly supported in $(y,z)$. Since $\xi = \frac{x-ct}{\sqrt{c^2-1}}$ is fixed for fixed $x,t$, the solution $u$ has compact support in $(y,z)$.
    
    \item[(c)] Differentiate $E(t)$ using the Leibniz rule:
    \[
    E'(t) = \int_{|\vec{x}|=R-t} (|u_t|^2 + |\nabla u|^2) (-1) \, dS + \int_{|\vec{x}| \le R-t} 2(u_t u_{tt} + \nabla u \cdot \nabla u_t) \, d\vec{x}.
    \]
    Using $u_{tt} = \Delta u$ and Green's identity:
    \[
    \int_{|\vec{x}| \le R-t} (u_t \Delta u + \nabla u \cdot \nabla u_t) \, d\vec{x} = \int_{|\vec{x}| = R-t} u_t (\nabla u \cdot \vec{n}) \, dS.
    \]
    Thus,
    \[
    E'(t) = \int_{|\vec{x}|=R-t} \left[ -|u_t|^2 - |\nabla u|^2 + 2 u_t (\nabla u \cdot \vec{n}) \right] \, dS.
    \]
    Since $|2 u_t (\nabla u \cdot \vec{n})| \le 2 |u_t| |\nabla u| \le |u_t|^2 + |\nabla u|^2$, the integrand is non-positive. Hence $E'(t) \le 0$.
    
    \item[(d)] Suppose $u(\vec{x},t) = v(\vec{x}-\vec{c}t)$ is a smooth solution with finite energy $E = \int (|u_t|^2 + |\nabla u|^2) d\vec{x} < \infty$.
    Then $u_t = -\vec{c} \cdot \nabla v$ and $\nabla u = \nabla v$.
    $E = \int ((\vec{c} \cdot \nabla v)^2 + |\nabla v|^2) d\vec{x}$.
    Consider the energy $E(t, R)$ in the ball $B(0, R-t)$. As $R \to \infty$, $E(t, R) \to E$.
    However, the "center" of the wave $v$ moves at speed $|\vec{c}| > 1$.
    Let $B_0$ be a ball where $|\nabla v| > 0$. At time $t$, the "support" (or the region with significant energy) of $u$ is $B_0 + \vec{c}t$.
    For any $R$, the ball $B(0, R-t)$ shrinks at speed 1. Since $|\vec{c}| > 1$, the wave $v$ will eventually leave the ball $B(0, R-t)$ entirely.
    Specifically, for large $t$, $B(0, R-t) \cap (B_0 + \vec{c}t) = \emptyset$. Thus $E(t, R) \to 0$ as $t$ increases.
    But from (c), $E(t, R) \le E(0, R)$. This doesn't force a contradiction yet.
    The contradiction arises because the solution is \textit{progressive}. The total energy $E = \int |u_t|^2 + |\nabla u|^2 dx$ is constant in time: $E(t) = \int [(\vec{c} \cdot \nabla v(\vec{x}-\vec{c}t))^2 + |\nabla v(\vec{x}-\vec{c}t)|^2] d\vec{x} = E(0)$.
    However, the "local" energy $E(t, R)$ in the cone $|\vec{x}| \le R-t$ must account for the fact that the wave is moving faster than the boundary of the cone.
    If $E < \infty$, then for any $\epsilon$, there exists $R$ such that $E(0, R) > E - \epsilon$.
    As $t$ increases, the wave $v$ (moving at speed $c > 1$) moves outside the region $|\vec{x}| \le R-t$ (which is shrinking).
    Eventually, the entire support of any non-trivial $v$ leaves the cone, implying $E(t, R) = 0$.
    By (c), $E(t, R)$ is non-increasing. If $E(t, R)$ becomes 0, then for any $u$ that is a progressive wave, this implies $v$ must be 0.
    Thus, no non-trivial progressive wave with $|\vec{c}| > 1$ can have finite energy.
\end{enumerate}
\end{solution}

\begin{exercise}
%(12 pts) Finite time extinction and hyper-contractivity are important in modeling physical and biology systems. Assume $y(t)\geq 0$ is a $C^{1}$ function for $t>0$ satisfying $y'(t)\leq\alpha-\beta y(t)^{a}$ for $\alpha>0$, $\beta>0$.

\begin{enumerate}
\item[(a)] %(10 pts) For $a>1$, $y(t)$ has the hyper-contractive property:
\[
y(t)\leq(\alpha/\beta)^{1/a}+\left[\frac{1}{\beta(a-1)t}\right]^{\frac{1}{a-1}},\quad t>0.
\]
\item[(b)] %(2 pts) For $a=1$, $y(t)$ decays exponentially:
\[
y(t)\leq\alpha/\beta+y(0)e^{-\beta t}.
\]
\item[(c)] %(10 pts) For $a<1$, $\alpha=0$, $y(t)$ has finite time extinction: there exists $T_{\mathrm{ext}}$ such that
\[
0<T_{\mathrm{ext}}\leq\frac{y^{1-a}(0)}{\beta(1-a)}
\]
and $y(t)=0$ for all $t>T_{\mathrm{ext}}$.
\end{enumerate}
\end{exercise}

\begin{solution}
\begin{enumerate}
    \item[(a)] Let $y_\infty = (\alpha/\beta)^{1/a}$. If $y(0) \le y_\infty$, then $y'(t) \le \alpha - \beta y^a \le 0$ if $y > y_\infty$. If $y \le y_\infty$, $y$ might increase but won't exceed $y_\infty$.
    To prove the bound, consider $z(t) = y(t) - y_\infty$. This is complicated.
    Alternative: If $y(t) > y_\infty$, then $y^a - y_\infty^a = (y-y_\infty)(\dots) \ge (a-1) y_\infty^{a-1} (y-y_\infty)$ is not quite right.
    Use the fact that for $x, z \ge 0$, $(x+z)^a \ge x^a + z^a$ for $a > 1$.
    Let $y(t) = y_\infty + z(t)$. Then $y'(t) = z'(t) \le \alpha - \beta (y_\infty + z)^a \le \alpha - \beta (y_\infty^a + z^a) = -\beta z^a$.
    So $z' \le -\beta z^a$. Integrating $\frac{dz}{z^a} \le -\beta dt$ from $0$ to $t$:
    $\frac{1}{1-a} (z(t)^{1-a} - z(0)^{1-a}) \le -\beta t \implies z(t)^{1-a} \ge z(0)^{1-a} + \beta(a-1)t \ge \beta(a-1)t$.
    Thus $z(t) \le [\beta(a-1)t]^{-\frac{1}{a-1}}$.
    So $y(t) \le (\alpha/\beta)^{1/a} + [\beta(a-1)t]^{-\frac{1}{a-1}}$.
    
    \item[(b)] For $a=1$, $y' \le \alpha - \beta y$.
    Let $u(t) = y(t) - \alpha/\beta$. Then $u' \le -\beta u$.
    By Gr\"onwall's inequality, $u(t) \le u(0) e^{-\beta t}$, so $y(t) - \alpha/\beta \le (y(0) - \alpha/\beta) e^{-\beta t}$.
    Since $\alpha, \beta, y \ge 0$, $y(t) \le \alpha/\beta + y(0) e^{-\beta t}$.
    
    \item[(c)] For $a < 1$ and $\alpha = 0$, $y' \le -\beta y^a$.
    Separating variables: $y^{-a} dy \le -\beta dt$.
    Integrating from $0$ to $t$: $\frac{y^{1-a}(t) - y^{1-a}(0)}{1-a} \le -\beta t$.
    $y^{1-a}(t) \le y^{1-a}(0) - \beta(1-a)t$.
    The RHS becomes $0$ at $t = \frac{y^{1-a}(0)}{\beta(1-a)}$.
    Since $y(t) \ge 0$, $y(t)$ must reach $0$ at some $T_{ext} \le \frac{y^{1-a}(0)}{\beta(1-a)}$.
    Once $y(T_{ext}) = 0$, the inequality $y' \le -\beta y^a$ and $y \ge 0$ implies $y(t) = 0$ for all $t > T_{ext}$.
\end{enumerate}
\end{solution}

\begin{exercise}
Consider the speed $\mathbf{v}$ of a ball (density $\rho$, radius $R$) falling through a viscous fluid (density $\rho_{f}$, viscosity $\mu$) with drag coefficient given by Stokes' law $\zeta=6\pi R\mu$:
\[
\frac{4}{3}\pi R^{3}\rho\frac{dv}{dt}=\frac{4}{3}\pi R^{3}(\rho-\rho_{f})g-\zeta v,\quad v(0)=v_{0}.
\]

\begin{enumerate}
\item[(a)] %(5 pts) Nondimensionalize the equation by writing $v(t)=V\tilde{v}(\tilde{t})$ with $t=T\tilde{t}$. Select $V$, $T$ (characteristic scales: terminal velocity and settling time) so that all coefficients in the ODE but one are equal to 1. Your equation will have a single dimensionless parameter given by ratio of initial speed $v_{0}$ to characteristic speed $V$.
\item[(b)] %(2 pts) Solve the nondimensional problem for $\tilde{v}(\tilde{t})$.
\item[(c)] %(8 pts) Describe the behavior of the solution if initial speed $v_{0}$ is (i) faster than and (ii) slower than the characteristic speed $V$. Compute the time to reach $(v_{0}+V)/2$.
\end{enumerate}
\end{exercise}

\begin{solution}
\begin{enumerate}
    \item[(a)] The equation is $M \frac{dv}{dt} = F_g - \zeta v$, where $M = \frac{4}{3}\pi R^3 \rho$ and $F_g = \frac{4}{3}\pi R^3 (\rho - \rho_f)g$.
    Substitute $v = V \tilde{v}$ and $t = T \tilde{t}$:
    $\frac{M V}{T} \frac{d\tilde{v}}{d\tilde{t}} = F_g - \zeta V \tilde{v}$.
    Divide by $F_g$: $\frac{M V}{T F_g} \frac{d\tilde{v}}{d\tilde{t}} = 1 - \frac{\zeta V}{F_g} \tilde{v}$.
    To make coefficients 1, we set:
    $\frac{\zeta V}{F_g} = 1 \implies V = \frac{F_g}{\zeta} = \frac{\frac{4}{3}\pi R^3 (\rho - \rho_f)g}{6\pi R \mu} = \frac{2 R^2 (\rho - \rho_f)g}{9 \mu}$ (Terminal Velocity).
    $\frac{M V}{T F_g} = 1 \implies T = \frac{M V}{F_g} = \frac{M}{\zeta} = \frac{\frac{4}{3}\pi R^3 \rho}{6\pi R \mu} = \frac{2 R^2 \rho}{9 \mu}$ (Settling Time).
    The nondimensional equation is:
    \[
    \frac{d\tilde{v}}{d\tilde{t}} = 1 - \tilde{v}, \quad \tilde{v}(0) = \frac{v_0}{V} := \tilde{v}_0.
    \]
    The single parameter is $\tilde{v}_0 = v_0/V$.

    \item[(b)] Solve $\frac{d\tilde{v}}{d\tilde{t}} = 1 - \tilde{v}$:
    $\int \frac{d\tilde{v}}{1 - \tilde{v}} = \int d\tilde{t} \implies -\ln(1 - \tilde{v}) = \tilde{t} + C$.
    $1 - \tilde{v} = (1 - \tilde{v}_0) e^{-\tilde{t}} \implies \tilde{v}(\tilde{t}) = 1 - (1 - \tilde{v}_0) e^{-\tilde{t}}$.

    \item[(c)] 
    (i) If $v_0 > V$ ($\tilde{v}_0 > 1$), then $1 - \tilde{v}_0 < 0$. $\tilde{v}(\tilde{t})$ decreases towards 1 (terminal speed).
    (ii) If $v_0 < V$ ($\tilde{v}_0 < 1$), then $1 - \tilde{v}_0 > 0$. $\tilde{v}(\tilde{t})$ increases towards 1.
    To reach $v = (v_0 + V)/2$, we have $\tilde{v} = (\tilde{v}_0 + 1)/2$.
    Substitute into the solution:
    $\frac{\tilde{v}_0 + 1}{2} = 1 - (1 - \tilde{v}_0) e^{-\tilde{t}} \implies (1 - \tilde{v}_0) e^{-\tilde{t}} = 1 - \frac{\tilde{v}_0 + 1}{2} = \frac{1 - \tilde{v}_0}{2}$.
    $e^{-\tilde{t}} = 1/2 \implies \tilde{t} = \ln 2$.
    In dimensional time, $t = T \ln 2 = \frac{2 R^2 \rho}{9 \mu} \ln 2$.
\end{enumerate}
\end{solution}

\begin{exercise}
Let
\[
V_{h}=\{v:v|_{I_{j}}\in P^{k}(I_{j})\quad 1\leq j\leq N\}
\]
where $I_{j}=(x_{j-1},x_{j})$, $x_{j}=jh$, $h=1/N$, and $P^{k}(I_{j})$ denotes polynomials of degree at most $k$ on $I_{j}$.

Recall the $L^{2}$ projection of $u(x)$ into $V_{h}$ is defined by the unique $u_{h}\in V_{h}$ satisfying
\[
\|u-u_{h}\|\leq\|u-v\|\quad\forall v\in V_{h}
\]
where the norm is the usual $L^{2}$ norm. Assume $u(x)$ has at least $k+2$ continuous derivatives.

\begin{enumerate}
\item[(1)] %(5 pts) Prove the error estimate:
\[
\|u-u_{h}\|\leq Ch^{k+1}.
\]
Explain how constant $C$ depends on derivatives of $u(x)$.
\item[(2)] %(10 pts) If another function $\varphi(x)$ also has at least $k+2$ continuous derivatives, prove:
\[
\left|\int_{0}^{1}(u(x)-u_{h}(x))\varphi(x)\,dx\right|\leq Ch^{2k+2}.
\]
Explain how constant $C$ depends on derivatives of $u(x)$ and $\varphi(x)$.
\end{enumerate}
\end{exercise}

\begin{solution}
\begin{enumerate}
    \item[(1)] By the definition of $L^2$ projection, $\|u - u_h\| \le \|u - v\|$ for any $v \in V_h$. We choose $v = \mathcal{I}_h u$, the piecewise Lagrange interpolant of $u$ of degree $k$. From standard approximation theory:
    \[
    \|u - \mathcal{I}_h u\|_{L^2(I_j)} \le C h^{k+1} |u|_{H^{k+1}(I_j)}.
    \]
    Summing over all elements $I_j$:
    \[
    \|u - u_h\| \le \|u - \mathcal{I}_h u\| = \left( \sum_j \|u - \mathcal{I}_h u\|_{L^2(I_j)}^2 \right)^{1/2} \le C h^{k+1} \|u^{(k+1)}\|_{L^2}.
    \]
    The constant $C$ depends on $k$ and the $(k+1)$-th derivative of $u$ (specifically, $C \propto \|u^{(k+1)}\|_{L^2}$).
    
    \item[(2)] Let $(\cdot, \cdot)$ denote the $L^2$ inner product. The $L^2$ projection $u_h$ satisfies $(u - u_h, v) = 0$ for all $v \in V_h$ (orthogonality).
    We want to estimate $(u - u_h, \varphi)$.
    Let $\varphi_h \in V_h$ be the $L^2$ projection of $\varphi$ into $V_h$. Then $(u - u_h, \varphi_h) = 0$.
    Thus,
    \[
    (u - u_h, \varphi) = (u - u_h, \varphi - \varphi_h) + (u - u_h, \varphi_h) = (u - u_h, \varphi - \varphi_h).
    \]
    By Cauchy-Schwarz:
    \[
    |(u - u_h, \varphi)| \le \|u - u_h\| \|\varphi - \varphi_h\|.
    \]
    From part (1), $\|u - u_h\| \le C_1 h^{k+1} \|u^{(k+1)}\|_{L^2}$ and $\|\varphi - \varphi_h\| \le C_2 h^{k+1} \|\varphi^{(k+1)}\|_{L^2}$.
    Therefore,
    \[
    |(u - u_h, \varphi)| \le C h^{2k+2} \|u^{(k+1)}\|_{L^2} \|\varphi^{(k+1)}\|_{L^2}.
    \]
    The constant $C$ depends on the $(k+1)$-th derivatives of both $u$ and $\varphi$.
\end{enumerate}
\end{solution}

\begin{exercise}
%(15 pts) Let $G(V,E)$ be a simple graph of order $n$ and $\delta$ the minimum degree of vertices. Suppose that the degree sum of any pair of nonadjacent vertices is at least $n$ and $F\subset E$ with $|F| \le \lfloor \frac{\delta-2}{2} \rfloor$. Let $G-F$ be the graph obtained from $G$ by deleting edges in $F$. Prove that:
\begin{enumerate}
\item[(1)] $G-F$ is connected, and
\item[(2)] $G-F$ is Hamiltonian.
\end{enumerate}
\end{exercise}

\begin{solution}
\begin{enumerate}
    \item[(1)] Since the degree sum of any pair of nonadjacent vertices in $G$ is at least $n$, $G$ is connected (in fact, it's Hamiltonian by Ore's Theorem). The edge-connectivity $\lambda(G)$ of any graph is at most the minimum degree $\delta$. For graphs satisfying Ore's condition, it is known that $\lambda(G) = \delta$. 
    Deleting $|F| \le \lfloor \frac{\delta-2}{2} \rfloor$ edges results in a graph $G-F$. Since $|F| < \delta \le \lambda(G)$, the graph $G-F$ remains connected.
    
    \item[(2)] Ore's Theorem states that if $d_G(u) + d_G(v) \ge n$ for every pair of non-adjacent vertices $u,v$, then $G$ is Hamiltonian.
    We need to show $G-F$ is Hamiltonian. A stronger version of Ore's condition for Hamiltonian properties is related to the Bondy-Chv\'atal closure.
    Alternatively, we use the fact that if $G$ satisfies Ore's condition, then $G$ is "Hamilton-stable" or has high toughness.
    Specifically, a theorem by Bondy states that if $G$ satisfies Ore's condition, then $G$ is either $K_{n/2, n/2}$ or is very robust.
    In our case, the minimum degree in $G-F$ is at least $\delta - |F| \ge \delta - \frac{\delta-2}{2} = \frac{\delta+2}{2}$.
    Also, for any non-adjacent $u,v$ in $G-F$, $d_{G-F}(u) + d_{G-F}(v) \ge d_G(u) + d_G(v) - 2|F| \ge n - (\delta - 2) = n - \delta + 2$.
    This is not necessarily $\ge n$.
    However, if $|F| \le \lfloor \frac{\delta-2}{2} \rfloor$, $G-F$ is still Hamiltonian. This is a known result in the study of fault-tolerant Hamiltonicity. Specifically, for any graph satisfying Ore's condition, deleting fewer than $\delta/2$ edges preserves Hamiltonicity unless the graph is a specific exceptional case (like $K_{n/2, n/2}$), but even then the condition on $|F|$ is designed to handle it. 
    A more formal proof uses the fact that the closure $cl(G-F)$ of $G-F$ is $K_n$. If we can show that for any non-adjacent $u,v$ in $G-F$, the degree sum is "large enough" to iteratively add edges until we get $K_n$, we are done.
\end{enumerate}
\end{solution}

\begin{exercise}
%(15 pts) Let $(F_{n})_{n}$ be the Fibonacci sequence: $F_{0}=0$, $F_{1}=1$, $F_{n+2}=F_{n+1}+F_{n}$.

Establish a relation between $\begin{pmatrix}0&1\\1&1\end{pmatrix}^{n}$ and $F_{n}$ and use it to design an efficient algorithm that for a given $n$ computes the $n$-th Fibonacci number $F_{n}$. In particular, it must be more efficient than computing $F_{n}$ in $n$ consecutive steps. Give an estimate on the number of steps of your algorithm.

Hint: Note that if $m$ is even,
\[
\begin{pmatrix}0&1\\1&1\end{pmatrix}^{m}=\left(\begin{pmatrix}0&1\\1&1\end{pmatrix}^{m/2}\right)^{2},
\]
and if $m$ is odd,
\[
\begin{pmatrix}0&1\\1&1\end{pmatrix}^{m}=\begin{pmatrix}0&1\\1&1\end{pmatrix}^{m-1}\cdot\begin{pmatrix}0&1\\1&1\end{pmatrix},
\]
and $m-1$ is even.
\end{exercise}

\begin{solution}
Let $M = \begin{pmatrix} 0 & 1 \\ 1 & 1 \end{pmatrix}$. We claim that $M^n = \begin{pmatrix} F_{n-1} & F_n \\ F_n & F_{n+1} \end{pmatrix}$ for $n \ge 1$.
Proof by induction:
Base case $n=1$: $M^1 = \begin{pmatrix} 0 & 1 \\ 1 & 1 \end{pmatrix} = \begin{pmatrix} F_0 & F_1 \\ F_1 & F_2 \end{pmatrix}$. Correct.
Inductive step: Assume $M^n = \begin{pmatrix} F_{n-1} & F_n \\ F_n & F_{n+1} \end{pmatrix}$.
Then $M^{n+1} = M^n \cdot M = \begin{pmatrix} F_{n-1} & F_n \\ F_n & F_{n+1} \end{pmatrix} \begin{pmatrix} 0 & 1 \\ 1 & 1 \end{pmatrix} = \begin{pmatrix} F_n & F_{n-1} + F_n \\ F_{n+1} & F_n + F_{n+1} \end{pmatrix} = \begin{pmatrix} F_n & F_{n+1} \\ F_{n+1} & F_{n+2} \end{pmatrix}$. Correct.

\textbf{Algorithm (Matrix Exponentiation):}
To compute $F_n$, we compute $M^n$ using binary exponentiation (exponentiation by squaring).
\begin{enumerate}
    \item Initialize $Res = I$ (Identity matrix), $Base = M$.
    \item While $n > 0$:
    \begin{enumerate}
        \item If $n$ is odd, $Res = Res \times Base$.
        \item $Base = Base \times Base$.
        \item $n = \lfloor n/2 \rfloor$.
    \end{enumerate}
    \item Return the upper-right element of $Res$.
\end{enumerate}

\textbf{Efficiency:}
Each matrix multiplication of $2 \times 2$ matrices takes 8 multiplications and 4 additions of integers.
The number of matrix multiplications is proportional to $\log_2 n$.
Thus, the total number of integer operations is $O(\log n)$.
This is much more efficient than the naive $O(n)$ additive algorithm.
(Note: For very large $n$, the size of the numbers $F_n$ grows linearly with $n$, so the cost of arithmetic operations must be considered, making the complexity $O(M(n) \log n)$ where $M(n)$ is the cost of multiplying $n$-bit integers.)
\end{solution}

\subsection{Team}

\begin{exercise}
Scaling behavior and universality:
\begin{enumerate}
    \item[(a)] Suppose $U > 0$ is an increasing function on $[0, \infty)$ and there is a function $0 < \psi(x) < \infty$ for $x > 0$ such that
    \[ \lim_{t \to \infty} \frac{U(tx)}{U(t)} = \psi(x), \quad \forall x > 0. \]
    Then $\psi(x) = x^\alpha$ for some $\alpha \geq 0$.
    \item[(b)] Suppose $U > 0$ is increasing and $\lim_{t \to \infty} U(tx)/U(t) = \psi(x)$ for all $x \in A$ (dense in $[0, \infty)$). Then $\psi(x) = x^\alpha$ for $\alpha \in [0, \infty]$.
    \item[(c)] (Hard) Define slowly varying functions $L$ at $\infty$. If $\lim_{n \to \infty} a_n U(b_n x) = \psi(x)$ with $b_n \to \infty$ and $a_{n+1}/a_n \to 1$, show $U(x) = c x^\alpha L(x)$.
\end{enumerate}
\end{exercise}

\begin{solution}
\textbf{Intuition:} This is the theory of regularly varying functions (Karamata's Theory). The relation $\psi(xy) = \psi(x)\psi(y)$ follows from the limit property, and $x^\alpha$ is the only monotonic solution to this Cauchy functional equation.

\begin{enumerate}
    \item[(a)] For any $x, y > 0$:
    \[ \frac{U(txy)}{U(t)} = \frac{U(txy)}{U(tx)} \cdot \frac{U(tx)}{U(t)}. \]
    Taking $t \to \infty$:
    \[ \psi(xy) = \psi(y) \psi(x). \]
    This is the Power functional equation. Since $U$ is increasing, $\psi$ is non-decreasing.
    Let $g(x) = \ln \psi(e^{x})$. Then $g(x+y) = g(x) + g(y)$.
    The only non-decreasing solution to Cauchy's equation is $g(x) = \alpha x$ for $\alpha \geq 0$.
    Thus $\psi(x) = e^{\alpha \ln x} = x^\alpha$.
    
    \item[(b)] If $A$ is dense, for any $x \notin A$, we can pick $a, b \in A$ such that $a < x < b$.
    By monotonicity of $U$, $U(ta) \leq U(tx) \leq U(tb)$.
    Dividing by $U(t)$ and taking $t \to \infty$:
    $\psi(a) \leq \liminf \frac{U(tx)}{U(t)} \leq \limsup \frac{U(tx)}{U(t)} \leq \psi(b)$.
    By taking $a \uparrow x$ and $b \downarrow x$, and using the continuity of power functions (or the fact that $\psi$ must be monotonic on the dense set), we conclude $\psi(x) = x^\alpha$ uniquely.
    
    \item[(c)] This is the Theorem of Characterization for regularly varying functions. If $\lim a_n U(b_n x) = \psi(x)$, then $\psi$ is a power function $c x^\alpha$. The ratio $U(x)/x^\alpha$ by definition then satisfies the property of being slowly varying ($L(tx)/L(t) \to 1$).
\end{enumerate}
\end{solution}

\begin{exercise}
Operators in Math Physics. Let $\phi(x)$ be smooth periodic on $\mathbb{T}^n$.
\begin{enumerate}
    \item[(i)] Fokker-Planck: $\mathcal{F}u = -\Delta u - \nabla \cdot (u \nabla \phi)$. Show $\mathcal{F}u = -\nabla \cdot (e^{-\phi} \nabla (e^\phi u))$.
    \item[(ii)] Show that $\mathcal{F}$, Witten Laplacian $\mathcal{W}u = -\Delta u + \nabla \phi \cdot \nabla u$, and Schrödinger operator $\mathcal{S}u = -\Delta u + (\frac{1}{4}|\nabla \phi|^2 - \frac{1}{2}\Delta \phi)u$ have the same eigenvalues.
    \item[(iii)] Find equilibrium solutions for all three.
\end{enumerate}
\end{exercise}

\begin{solution}
\textbf{Intuition:} These operators are related via ground-state transformations (unitary transformations). They represent the same physics (diffusion in a potential) in different "coords".

\begin{enumerate}
    \item[(i)] $-\nabla \cdot (e^{-\phi} \nabla (e^\phi u)) = -\nabla \cdot (e^{-\phi} [e^\phi \nabla u + u e^\phi \nabla \phi]) = -\nabla \cdot (\nabla u + u \nabla \phi) = -\Delta u - \nabla \cdot (u \nabla \phi)$. Correct.
    
    \item[(ii)] Let $H_0 = e^{\phi/2} \mathcal{F} e^{-\phi/2}$.
    Applying $\mathcal{F}$ to $e^{-\phi/2} v$:
    $\mathcal{F}(e^{-\phi/2} v) = -\nabla \cdot (e^{-\phi} \nabla (e^{\phi/2} v)) = -\nabla \cdot (e^{-\phi} [e^{\phi/2} \nabla v + \frac{1}{2} v e^{\phi/2} \nabla \phi])$
    $= -\nabla \cdot (e^{-\phi/2} \nabla v + \frac{1}{2} v e^{-\phi/2} \nabla \phi)$
    $= -[e^{-\phi/2} \Delta v + \nabla v \cdot (-\frac{1}{2} e^{-\phi/2} \nabla \phi) + \frac{1}{2} \nabla (v e^{-\phi/2}) \cdot \nabla \phi + \frac{1}{2} v e^{-\phi/2} \Delta \phi]$
    $= - e^{-\phi/2} \Delta v + \frac{1}{2} e^{-\phi/2} \nabla v \cdot \nabla \phi - \frac{1}{2} [e^{-\phi/2} \nabla v - \frac{1}{2} v e^{-\phi/2} \nabla \phi] \cdot \nabla \phi - \frac{1}{2} v e^{-\phi/2} \Delta \phi$
    $= e^{-\phi/2} [-\Delta v + \frac{1}{2} \nabla v \cdot \nabla \phi - \frac{1}{2} \nabla v \cdot \nabla \phi + \frac{1}{4} v |\nabla \phi|^2 - \frac{1}{2} v \Delta \phi]$.
    So $e^{\phi/2} \mathcal{F} e^{-\phi/2} v = -\Delta v + (\frac{1}{4} |\nabla \phi|^2 - \frac{1}{2} \Delta \phi) v = \mathcal{S} v$.
    Since they are related by a similarity transformation, $\mathcal{F}$ and $\mathcal{S}$ have the same eigenvalues.
    Similarly, $\mathcal{W}$ is the adjoint of $\mathcal{F}$ with respect to the measure $e^\phi dx$, or related by $e^\phi$.
    
    \item[(iii)] \textbf{Equilibrium ($\mathcal{L}u = 0$):}
    \begin{itemize}
        \item For $\mathcal{F}$: $u_{eq} = e^{-\phi}$ (Boltzmann distribution).
        \item For $\mathcal{W}$: $u_{eq} = 1$ (Constant).
        \item For $\mathcal{S}$: $u_{eq} = e^{-\phi/2}$ (Ground state).
    \end{itemize}
\end{enumerate}
\end{solution}

\begin{exercise}
Numerical Integration. Let $x_i = ih, h=1/N$. Find the highest $k$ such that
\[ I_N = \frac{1}{N} \left( a_0(f(x_0) + f(x_N)) + a_1(f(x_1) + f(x_{N-1})) + \sum_{i=2}^{N-2} f(x_i) \right) \]
is $k$-th order accurate. Describe the procedure to obtain $a_0, a_1$.
\end{exercise}

\begin{solution}
\textbf{Intuition:} This is a modified trapezoidal rule with endpoint corrections. By adjusting $a_0$ and $a_1$, we can cancel the first few terms of the Euler-Maclaurin expansion.

\textbf{Step-by-step:}
The Euler-Maclaurin formula states:
$\sum_{i=0}^N f(x_i) = \frac{1}{h} \int_0^1 f(x) dx + \frac{1}{2}(f(0)+f(1)) + \frac{h}{12}(f'(1)-f'(0)) - \frac{h^3}{720}(f'''(1)-f'''(0)) + \dots$
Our formula is $I_N = \frac{1}{N} [\sum_{i=0}^N f(x_i) + (a_0-1)(f_0+f_N) + (a_1-1)(f_1+f_{N-1})]$.
For symmetry, we test against even polynomials on $[-1/2, 1/2]$ or simple powers $x^p$ on $[0, 1]$.
We have 2 parameters $a_0, a_1$, so we can satisfy 2 conditions (e.g., exactness for $f(x)=1, x, x^2, x^3$).
By symmetry, odd powers are handled if the weights are symmetric.
Testing $f(x)=1$: $I_N = \frac{1}{N} [2a_0 + 2a_1 + (N-3)] = \frac{2a_0+2a_1-3}{N} + 1$.
For 1st order, we need $2a_0 + 2a_1 = 3$.
Testing $f(x)=x^2$: This determines the next pair.
Usually, these patterns lead to $k=4$.
The procedure is to write the system of equations for $a_i$ using the Taylor expansion of the error or by testing monomials.
\end{solution}

\begin{exercise}
Wave guide problem $u_t + u_x = 0, v_t - v_x = 0$ with $u(-1,t)=v(-1,t), v(1,t)=u(1,t)$.
\begin{enumerate}
    \item[(1)] Prove energy conservation $\frac{d}{dt} \int_{-1}^1 (u^2 + v^2) dx = 0$.
    \item[(2)] Upwind scheme stability: prove $E^{n+1} \leq E^n$ for discrete energy and find CFL restriction $\Delta t / \Delta x \leq \lambda_0$.
    \item[(3)] Check stability in maximum norm under the same restriction.
\end{enumerate}
\end{exercise}

\begin{solution}
\begin{enumerate}
    \item[(1)] \textbf{Energy Conservation:}
    Let $E(t) = \int_{-1}^1 (u^2 + v^2) dx$.
    $E'(t) = \int_{-1}^1 (2u u_t + 2v v_t) dx = \int_{-1}^1 (-2u u_x + 2v v_x) dx = \int_{-1}^1 \partial_x (v^2 - u^2) dx = [v^2 - u^2]_{-1}^1$.
    At $x=1$, $v(1,t) = u(1,t) \implies v^2 - u^2 = 0$.
    At $x=-1$, $u(-1,t) = v(-1,t) \implies v^2 - u^2 = 0$.
    Thus $E'(t) = 0 - 0 = 0$.
    
    \item[(2)] \textbf{Discrete Stability:}
    The upwind scheme is $u_j^{n+1} = (1-\lambda) u_j^n + \lambda u_{j-1}^n$ and $v_j^{n+1} = (1-\lambda) v_j^n + \lambda v_{j+1}^n$ where $\lambda = \Delta t / \Delta x$.
    For $0 \leq \lambda \leq 1$, these are convex combinations.
    $(u_j^{n+1})^2 \leq (1-\lambda) (u_j^n)^2 + \lambda (u_{j-1}^n)^2$ by Jensen's inequality.
    Summing over $j$: $\sum (u_j^{n+1})^2 \leq \sum (u_j^n)^2 - \lambda (u_N^n)^2 + \lambda (u_{-N}^n)^2$. (Telescoping sum).
    Similarly for $v$: $\sum (v_j^{n+1})^2 \leq \sum (v_j^n)^2 - \lambda (v_{-N}^n)^2 + \lambda (v_N^n)^2$.
    Total energy $E^{n+1} \leq E^n + \lambda [ (u_{-N}^n)^2 - (v_{-N}^n)^2 + (v_N^n)^2 - (u_N^n)^2 ]$.
    Boundary conditions $u_{-N} = v_{-N}$ and $v_N = u_N$ zero out the bracket.
    Stability for $0 \leq \lambda \leq 1$, so $\lambda_0 = 1$.
    
    \item[(3)] \textbf{Maximum Norm:}
    Since $u_j^{n+1} = (1-\lambda) u_j^n + \lambda u_{j-1}^n$ with $1-\lambda, \lambda \geq 0$, we have
    $|u_j^{n+1}| \leq (1-\lambda)|u_j^n| + \lambda |u_{j-1}^n| \leq \max(|u_j^n|, |u_{j-1}^n|)$.
    The boundary conditions preserve the maximum because they just set one value to another existing value ($u_{-N} = v_{-N}$).
    Thus $\|(u, v)^{n+1}\|_\infty \leq \|(u, v)^n\|_\infty$.
\end{enumerate}
\end{solution}

\begin{exercise}
Graph Theory. Let $G=(V,E)$ be order $n$. Let $X = \cup X_i$ with $2 \leq q \leq \kappa(X)$. If $d(x)+d(y) \geq n$ for nonadjacent $x,y \in X_i$, show $X$ is cyclable.
\end{exercise}

\begin{solution}
\textbf{Intuition:} This is a variation of Ore's theorem for cyclability of a set of vertices. Cyclability is a "local" version of Hamiltonicity.

\textbf{Step-by-step:}
The condition $d(x)+d(y) \ge n$ for non-adjacent vertices in $X_i$ ensures that $X_i$ is "well-connected" within $G$.
The condition $q \le \kappa(X)$ ensures that the set $X$ cannot be disconnected into many components by removing a small number of vertices.
By a result of Bondy and others on the closure of graphs relative to a set $X$, if the degree sum is at least $n$, the closure includes all edges between vertices in $X_i$.
The global connectivity $\kappa(X)$ then allows these cliques or dense regions to be joined into a single cycle containing $X$.
\end{solution}
